\makeatletter
\@ifundefined{keepfig}{\newenvironment{keepfig}{\par\edef\keepfig@pi{\the\parindent}\noindent
  \begin{minipage}{\linewidth}\parindent=\keepfig@pi\relax}{\end{minipage}\par}}{}
\makeatother
\section{Многочлены степени 4}

Пусть $f\in F[X]$~--- сепарабельный многочлен степени $4$, $E$~--- его поле разложения, $x_1,x_2,x_3,x_4$~--- корни $f$ в $E$, и
$G=G_f=\Gal(E/F)$, которую мы отождествляем с подгруппой $\Sym(\{x_1,x_2,x_3,x_4\})=S_4$. Всюду в этом разделе
$\operatorname{char}F\ne2$ (иначе дискриминант бесполезен, \rf{4.4}). Главная идея: у $S_4$ есть нормальная подгруппа
\[
V=\{1,\ (1\,2)(3\,4),\ (1\,3)(2\,4),\ (1\,4)(2\,3)\},
\]
и факторгруппа $S_4/V\cong S_3$ <<видна>> через кубический многочлен~--- кубическую резольвенту. Тем самым задача для степени
$4$ сводится к степени $3$.

\subsection*{Три разбиения на пары}

Четыре корня разбиваются на две пары тремя способами: $\{12|34\}$, $\{13|24\}$, $\{14|23\}$. Сопоставим им элементы
\[
\theta_1=x_1x_2+x_3x_4,\qquad \theta_2=x_1x_3+x_2x_4,\qquad \theta_3=x_1x_4+x_2x_3 .
\]

\begin{proposition}[4.7j]
\begin{enumerate}[label=\textup{(\alph*)}]
\item $\theta_1-\theta_2=(x_1-x_4)(x_2-x_3)$, $\theta_1-\theta_3=(x_1-x_3)(x_2-x_4)$, $\theta_2-\theta_3=(x_1-x_2)(x_3-x_4)$; в частности,
$\theta_1,\theta_2,\theta_3$ попарно различны, и
\[
(\theta_1-\theta_2)(\theta_1-\theta_3)(\theta_2-\theta_3)=\prod_{1\le i<j\le4}(x_i-x_j)=\Delta(f).
\]
\item Каждая перестановка $\sigma\in S_4$ переставляет три разбиения; получаем сюръективный гомоморфизм $\pi\colon S_4\to S_3$
с ядром $V$, и $\sigma(\theta_i)=\theta_{\pi(\sigma)(i)}$.
\item Стабилизатор разбиения $\{1j|kl\}$~--- подгруппа $D^{(j-1)}$ порядка $8$:
\begin{align*}
D^{(1)}&=\{1,(1\,2),(3\,4),(1\,2)(3\,4),(1\,3)(2\,4),(1\,4)(2\,3),(1\,3\,2\,4),(1\,4\,2\,3)\}=\langle(1\,3\,2\,4),(1\,2)\rangle,\\
D^{(2)}&=\langle(1\,2\,3\,4),(1\,3)\rangle,\qquad D^{(3)}=\langle(1\,2\,4\,3),(1\,4)\rangle .
\end{align*}
Это все три силовские $2$-подгруппы $S_4$; каждая изоморфна диэдральной группе $D_4$ (группе симметрий квадрата), они
сопряжены, и $V=D^{(1)}\cap D^{(2)}\cap D^{(3)}$.
\end{enumerate}
\end{proposition}

\begin{proof}
(a) Например, $\theta_1-\theta_2=x_1x_2+x_3x_4-x_1x_3-x_2x_4=x_1(x_2-x_3)-x_4(x_2-x_3)$. Перемножив три равенства, получаем
произведение всех шести разностей $x_i-x_j$, $i<j$, причём каждая входит со знаком <<$+$>>.

(b) Перестановка корней переводит разбиение на пары в разбиение на пары, и это действие согласовано с композицией,
т.\,е.\ $\pi$~--- гомоморфизм. Транспозиция $(1\,2)$ сохраняет $\{12|34\}$ и меняет местами $\{13|24\}$ и $\{14|23\}$, а $(1\,3)$
сохраняет $\{13|24\}$ и меняет местами два других; две разные транспозиции порождают $S_3$, так что $\pi$ сюръективен.
Его ядро имеет порядок $24/6=4$ и содержит $V$ (каждое произведение двух независимых транспозиций сохраняет все три
разбиения), т.\,е.\ равно $V$. Равенство $\sigma(\theta_i)=\theta_{\pi(\sigma)(i)}$ очевидно из определения $\theta_i$.

(c) Стабилизатор точки при транзитивном действии на трёхэлементном множестве имеет индекс $3$, т.\,е.\ порядок $8$;
непосредственно проверяется, что перечисленные $8$ перестановок сохраняют $\{12|34\}$. Порядок $8$~--- максимальная степень
двойки, делящая $24$, так что это силовские $2$-подгруппы; по теореме Силова их число нечётно и делит $3$. Три
стабилизатора различны (например, $(1\,2)\in D^{(1)}\smallsetminus D^{(2)}$) и сопряжены (сопряжены точки), значит, других силовских
$2$-подгрупп нет. Пересечение трёх стабилизаторов~--- это ядро действия, т.\,е.\ $V$. Наконец, $D^{(1)}$ переставляет вершины
квадрата, обходимые в порядке $1,3,2,4$ (рис.~ниже), и является группой его симметрий.
\end{proof}

\begin{keepfig}
\begin{center}
\begin{tikzpicture}[every node/.style={font=\small}]
\foreach \dx/\pa/\pb/\pc/\pd/\lab/\thx in {0/1/3/2/4/{\{12|34\}}/{\theta_1},4.5/1/2/3/4/{\{13|24\}}/{\theta_2},9/1/2/4/3/{\{14|23\}}/{\theta_3}}{
\begin{scope}[xshift=\dx cm]
\draw[cold,thick] (0,0)--(1.6,0)--(1.6,1.6)--(0,1.6)--cycle;
\draw[cnew,dashed] (0,0)--(1.6,1.6); \draw[cnew,dashed] (1.6,0)--(0,1.6);
\node[below left] at (0,0) {$\pa$}; \node[below right] at (1.6,0) {$\pb$};
\node[above right] at (1.6,1.6) {$\pc$}; \node[above left] at (0,1.6) {$\pd$};
\node at (0.8,-0.75) {диагонали $\lab$, $\thx$};
\end{scope}}
\end{tikzpicture}
\end{center}
\figcap{Три силовские $2$-подгруппы $D^{(1)},D^{(2)},D^{(3)}$ группы $S_4$~--- группы симметрий трёх квадратов с вершинами $1,2,3,4$;
пары противоположных вершин (диагонали) образуют соответствующее разбиение на пары. Сопряжение $\tau D^{(i)}\tau^{-1}$ с
$\tau\in S_4$~--- это \emph{перенумерация} вершин $k\mapsto\tau(k)$, а не движение плоскости: оно переводит квадрат с диагоналями
$\{ij|kl\}$ в квадрат с диагоналями $\{\tau i\,\tau j|\tau k\,\tau l\}$, т.\,е.\ $D^{(i)}$ в $D^{(\pi(\tau)(i))}$. Например, $\tau=(1\,2\,3\,4)$ переводит
квадрат $1,3,2,4$ в квадрат $2,4,3,1$~--- это квадрат с диагоналями $\{14|23\}$, а не исходный. Подгруппа $V$, общая для трёх
квадратов,~--- группа симметрий \emph{прямоугольника} (не квадрата) с вершинами $1,3,2,4$: поворот на $180^\circ$~--- это
$(1\,2)(3\,4)$, отражения относительно средних линий~--- $(1\,3)(2\,4)$ и $(1\,4)(2\,3)$.}
\end{keepfig}

\begin{lemma}[4.8]
Поле $F[\theta_1,\theta_2,\theta_3]$~--- неподвижное поле подгруппы $G\cap V$. Следовательно, $M=F[\theta_1,\theta_2,\theta_3]$~--- расширение
Галуа поля $F$ с группой $G/(G\cap V)$:
\[
\begin{tikzcd}[row sep=small,column sep=large]
F \arrow[r,dash,"G/(G\cap V)"] & M=F[\theta_1,\theta_2,\theta_3] \arrow[r,dash,"G\cap V"] & E .
\end{tikzcd}
\]
\end{lemma}

\begin{proof}
По \rf{4.7j}(b) элемент $\sigma\in G$ действует на $\theta_1,\theta_2,\theta_3$ перестановкой $\pi(\sigma)$, а так как $\theta_i$ различны,
$\sigma$ оставляет на месте все $\theta_i$ тогда и только тогда, когда $\pi(\sigma)=1$, т.\,е.\ $\sigma\in V$. Значит, $\Gal(E/M)=G\cap V$, и
$M=E^{G\cap V}$ по основной теореме (\rf{3.17}). Подгруппа $G\cap V$ нормальна в $G$ (так как $V$ нормальна в $S_4$), поэтому
$M/F$~--- расширение Галуа с группой $G/(G\cap V)$ (\rf{3.17}).
\end{proof}

\begin{definition}[4.8a]
Многочлен $g(X)=(X-\theta_1)(X-\theta_2)(X-\theta_3)$ называется \emph{кубической резольвентой} многочлена $f$. Любой $\sigma\in G$
переставляет $\theta_i$ и потому сохраняет $g$; значит, коэффициенты $g$ лежат в $E^G=F$. Поле $M=F[\theta_1,\theta_2,\theta_3]$~---
поле разложения $g$ над $F$.
\end{definition}

Коэффициенты $g$ можно найти, раскрыв скобки и выразив симметрические функции от $x_i$ через коэффициенты $f$, но это
громоздко. Удобнее следующее рассуждение, которое заодно показывает, \emph{откуда} берётся резольвента: корень $\theta_1$
отвечает разложению $f$ на два квадратичных множителя $(X-x_1)(X-x_2)\cdot(X-x_3)(X-x_4)$.

\begin{lemma}[4.9]
Кубическая резольвента многочлена $f=X^4+bX^3+cX^2+dX+e$ равна
\[
g=X^3-cX^2+(bd-4e)X-b^2e+4ce-d^2 .
\]
Дискриминанты $f$ и $g$ равны, и даже $\Delta(f)=\Delta(g)$ при нумерации корней из \rf{4.7j}.
\end{lemma}

\begin{proof}
Положим $s=x_1+x_2$, $s'=x_3+x_4$, $P=x_1x_2$, $P'=x_3x_4$. Тогда $f=(X^2-sX+P)(X^2-s'X+P')$, и сравнение коэффициентов даёт
\[
s+s'=-b,\qquad ss'+P+P'=c,\qquad sP'+s'P=-d,\qquad PP'=e .
\]
Здесь $P+P'=\theta_1$, так что $ss'=c-\theta_1$. Кроме того, $(sP'+s'P)+(sP+s'P')=(s+s')(P+P')=-b\theta_1$, т.\,е.\ $sP+s'P'=d-b\theta_1$.
Перемножим:
\[
(sP'+s'P)(sP+s'P')=PP'(s^2+s'^2)+ss'(P^2+P'^2)=PP'\bigl((s+s')^2-2ss'\bigr)+ss'\bigl((P+P')^2-2PP'\bigr).
\]
Подставляя, получаем $-d(d-b\theta_1)=e\bigl(b^2-2c+2\theta_1\bigr)+(c-\theta_1)(\theta_1^2-2e)$, т.\,е.\ после раскрытия скобок
\[
\theta_1^3-c\theta_1^2+(bd-4e)\theta_1-b^2e+4ce-d^2=0 .
\]
То же верно для $\theta_2,\theta_3$ (перенумеруем корни). Итак, три различных (\rf{4.7j}(a)) элемента $\theta_i$~--- корни унитарного
кубического многочлена из формулировки, и он равен $g$. Наконец, по \rf{4.7j}(a)
$\Delta(g)=(\theta_1-\theta_2)(\theta_1-\theta_3)(\theta_2-\theta_3)=\Delta(f)$, и $D(g)=\Delta(g)^2=\Delta(f)^2=D(f)$.
\end{proof}

\begin{remark}[4.9a]
(a) Коэффициенты $g$~--- многочлены с целыми коэффициентами от $b,c,d,e$, так что формула \rf{4.9}~--- тождество, верное для
любого (не обязательно сепарабельного) $f$: достаточно применить её к <<общему>> многочлену $\prod(X-x_i)$ над
$\Q(x_1,\dots,x_4)$ с независимыми переменными $x_i$ (там $\theta_i$ различны) и затем подставить значения.

(b) \emph{Приведённая форма.} Замена $X\mapsto X-\frac b4$ (формула Тейлора в точке $-\frac b4$) переводит $f$ в
\[
X^4+pX^2+qX+r,\qquad p=c-\frac{3b^2}8,\quad q=d-\frac{bc}2+\frac{b^3}8,\quad r=e-\frac{bd}4+\frac{b^2c}{16}-\frac{3b^4}{256};
\]
поле разложения и группа Галуа при этом не меняются, а все $\theta_i$ сдвигаются на одно и то же число $-\frac{b^2}8$. Для
приведённого многочлена
\[
g=X^3-pX^2-4rX+(4pr-q^2).
\]

(c) \emph{Явная формула для дискриминанта.} Как и в (a), равенство $\Delta(f)=\Delta(g)$ из \rf{4.9}~--- тождество: для общего
многочлена $\prod(X-x_i)$ это равенство многочленов от $x_1,\dots,x_4$ с целыми коэффициентами, и подстановка корней
показывает, что $D(f)=D(g)$ для любого унитарного $f$ степени $4$ над любым полем, в любой характеристике.
Применяя к $g$ формулу для кубики (\rf{4.2a})
$D(X^3+AX^2+BX+C)=A^2B^2-4B^3-4A^3C-27C^2+18ABC$, для приведённого многочлена ($A=-p$, $B=-4r$, $C=4pr-q^2$) получаем
\[
D(X^4+pX^2+qX+r)=256r^3-128p^2r^2+144pq^2r-27q^4+16p^4r-4p^3q^2 ;
\]
при $q=0$ и при $p=0$ отсюда снова получаются формулы \rf{4.2f}(a),(b): $D(X^4+pX^2+r)=16r(p^4-8p^2r+16r^2)=16r(p^2-4r)^2$ и
$D(X^4+qX+r)=256r^3-27q^4$. Для общего многочлена $X^4+bX^3+cX^2+dX+e$ та же подстановка коэффициентов резольвенты из \rf{4.9}
($A=-c$, $B=bd-4e$, $C=-b^2e+4ce-d^2$) даёт формулу из $16$ слагаемых:
\begin{multline*}
D=256e^3-192bde^2-128c^2e^2+144cd^2e-27d^4+144b^2ce^2-6b^2d^2e-80bc^2de\\
+18bcd^3+16c^4e-4c^3d^2-27b^4e^2+18b^3cde-4b^3d^3-4b^2c^3e+b^2c^2d^2 .
\end{multline*}
\end{remark}

\subsection*{Транзитивные подгруппы $S_4$}

\begin{proposition}[4.9b]
Транзитивные подгруппы $S_4$~--- это $S_4$, $A_4$, $V$, три сопряжённые подгруппы $D^{(1)},D^{(2)},D^{(3)}\cong D_4$ (\rf{4.7j}(c)) и три
сопряжённые циклические подгруппы $\langle(1\,3\,2\,4)\rangle,\langle(1\,2\,3\,4)\rangle,\langle(1\,2\,4\,3)\rangle\cong C_4$. Для них:
\begin{center}
\begin{tabular}{@{}lcccl@{}}
\toprule
$G$ & $|G|$ & $G\cap V$ & $(G:G\cap V)$ & $G\subset A_4$?\\
\midrule
$S_4$ & $24$ & $V$ & $6$ & нет\\
$A_4$ & $12$ & $V$ & $3$ & да\\
$V$ & $4$ & $V$ & $1$ & да\\
$D_4$ & $8$ & $V$ & $2$ & нет\\
$C_4$ & $4$ & $C_2$ & $2$ & нет\\
\bottomrule
\end{tabular}
\end{center}
(для $C_4=\langle(1\,3\,2\,4)\rangle$ имеем $C_4\cap V=\langle(1\,2)(3\,4)\rangle$).
\end{proposition}

\begin{proof}
Если $G$ транзитивна, то $|G|=4\cdot|\Stab(1)|$ делится на $4$, т.\,е.\ $|G|\in\{4,8,12,24\}$.

$|G|=24$: $G=S_4$. $|G|=12$: подгруппа индекса $2$ в $S_4$ равна $A_4$ (это показано в \rf{3.27}: она нормальна, вместе с
каждым $\sigma$ содержит $\sigma^2$, а каждый $3$-цикл~--- квадрат $3$-цикла, и $3$-циклы порождают $A_4$; иначе: $S_4/G\cong C_2$
абелева, поэтому $G$ содержит коммутант $S_4$, равный $A_4$). Разбирать абстрактные группы порядка $12$ здесь не нужно;
они разобраны в \rf{4.9c}(d): все они, кроме $A_4$, содержат элемент порядка $6$, которого в $S_4$ нет. (Заметим ещё, что абстрактная группа $G$ порядка $12$ с четырьмя силовскими $3$-подгруппами
$P_1,\dots,P_4$ вкладывается в $S_4$ своим действием сопряжением на $\{P_1,\dots,P_4\}$: ядро лежит в $\bigcap N(P_i)=\bigcap P_i=1$
(так как $(G:N(P_i))=4$, т.\,е.\ $N(P_i)=P_i$), и образ порядка $12$~--- это $A_4$; значит, $G\cong A_4$.) $|G|=8$: $G$~---
силовская $2$-подгруппа, т.\,е.\ одна из $D^{(i)}$ (\rf{4.7j}(c)); все они транзитивны.

$|G|=4$: $G$~--- $2$-группа и по теореме Силова лежит в одной из $D^{(i)}$. Если $G$ циклическая, то её образующая имеет
порядок $4$, т.\,е.\ является $4$-циклом; три подгруппы, порождённые $4$-циклами (по два $4$-цикла в каждой), перечислены
выше, и они транзитивны. Если $G\cong C_2\times C_2$, то её неединичные элементы~--- транспозиции или произведения двух
независимых транспозиций. Если транспозиций в $G$ нет, то $G$ состоит из $1$ и трёх таких произведений, т.\,е.\ $G=V$ (она
транзитивна). Если $(i\,j)\in G$, то $G$ абелева и лежит в централизаторе $(i\,j)$, который имеет порядок $24/6=4$ и равен
$\{1,(i\,j),(k\,l),(i\,j)(k\,l)\}$ ($\{i,j,k,l\}=\{1,2,3,4\}$); эта группа не транзитивна (орбиты $\{i,j\}$ и $\{k,l\}$).

Столбцы таблицы проверяются непосредственно: $D^{(i)}\supset V$ (\rf{4.7j}(c)), $(1\,3\,2\,4)^2=(1\,2)(3\,4)$, а $4$-циклы и транспозиции
нечётны.
\end{proof}

\begin{remark}[4.9c: как искать транзитивные подгруппы $S_n$]
(a) Если $G\subset S_n$ транзитивна, то $n$ делит $|G|$ (формула орбит), а стабилизатор $H$ точки $1$ имеет индекс $n$.

(b) Обратно, абстрактная группа $G$ с подгруппой $H$ индекса $n$ действует левыми сдвигами на $n$-элементном множестве
$G/H$ транзитивно; ядро этого действия~--- $\bigcap_{g\in G}gHg^{-1}$ (наибольшая нормальная подгруппа, лежащая в $H$). Если оно
тривиально, получаем вложение $G\hookrightarrow S_n$ с транзитивным образом, и всякая транзитивная подгруппа $S_n$ получается
так (с $H=\Stab(1)$, причём выбор нумерации $G/H$ меняет образ на сопряжённый). Поэтому перечисление транзитивных
подгрупп $S_n$ с точностью до сопряжения сводится к перечислению пар <<группа $G$ порядка, делящегося на $n$; подгруппа
$H$ индекса $n$, не содержащая неединичных нормальных подгрупп $G$>> (с точностью до автоморфизмов $G$). Так сразу
отсекается, например, подгруппа $\{1,(1\,2),(3\,4),(1\,2)(3\,4)\}$ группы $S_4$: в абелевой группе все подгруппы нормальны,
так что для абелевой транзитивной $G$ стабилизатор точки тривиален и $|G|=n$, а здесь стабилизатор точки $1$ равен
$\{1,(3\,4)\}$.

(c) Полезные ограничения: при простом $n=p$ транзитивная группа содержит элемент порядка $p$, т.\,е.\ $p$-цикл (по теореме
Коши); силовские подгруппы и их число; знак перестановок. Для $n=5$ транзитивных подгрупп с точностью до сопряжения
тоже пять (\rf{4.18a}).

(d) \emph{Группы порядка $12$.} Пусть $|G|=12$, и $n_2,n_3$~--- числа силовских $2$- и $3$-подгрупп. По теореме Силова
$n_2\mid3$, $n_3\mid4$, $n_3\equiv1\pmod 3$, т.\,е.\ $n_2\in\{1,3\}$, $n_3\in\{1,4\}$. Случай $n_2=3$, $n_3=4$ невозможен: четыре подгруппы
порядка $3$ попарно пересекаются по $1$ и дают $8$ элементов порядка $3$, а каждая силовская $2$-подгруппа ($4$ элемента,
порядки которых делят $4$) содержится в множестве оставшихся $4$ элементов, т.\,е.\ совпадает с ним и потому единственна.
Значит, одна из силовских подгрупп, $N$, нормальна. Для силовской подгруппы $H$ другого порядка $N\cap H=1$, а $NH$~---
подгруппа порядка $12$, т.\,е.\ $G=N\rtimes_\psi H$, где $\psi\colon H\to\Aut N$~--- действие сопряжением. Напомним, что
$N\rtimes_\psi H\cong N\rtimes_{\psi\circ\alpha}H$ для любого $\alpha\in\Aut H$ (изоморфизм $(n,h)\mapsto(n,\alpha^{-1}(h))$).
\begin{itemize}
\item $n_2=n_3=1$: $G=N\times H$ абелева, $G\cong C_{12}$ или $C_2\times C_6$.
\item $n_2=1$, $n_3=4$: $H=\langle c\rangle\cong C_3$ действует на $N$ нетривиально (иначе $G=N\times H$, $H$ нормальна и $n_3=1$). Так
как $\Aut(C_4)\cong C_2$ не содержит элементов порядка $3$, имеем $N=\{1,a,b,ab\}\cong C_2\times C_2$. Автоморфизм $N$ переставляет
три неединичных элемента $a,b,ab$, и обратно, \emph{любая} их перестановка~--- автоморфизм (произведение двух различных из
них равно третьему, а квадрат каждого равен $1$); итак, $\Aut(C_2\times C_2)\cong S_3$. Элемент $\psi(c)$ порядка $3$~--- циклическая
перестановка: либо $cac^{-1}=b$, $cbc^{-1}=ab$ (цикл $a\mapsto b\mapsto ab\mapsto a$), либо $cac^{-1}=ab$, $cbc^{-1}=a$ (обратный цикл). Два варианта переходят друг в друга заменой образующей $c$ на $c^{-1}$
(автоморфизм $\alpha\colon c\mapsto c^{-1}$ группы $H$) и потому дают изоморфные группы. Это $A_4$ ($N=V$, $c=(1\,2\,3)$).
\item $n_2=3$, $n_3=1$: $N=\langle t\rangle\cong C_3$, и $\psi\colon H\to\Aut(C_3)=\{t\mapsto t^{\pm1}\}\cong C_2$ нетривиален (иначе $G$ абелева
и $n_2=1$). При $H=\langle s\rangle\cong C_4$ такой $\psi$ единствен ($sts^{-1}=t^{-1}$): $G\cong C_3\rtimes C_4$. При $H\cong C_2\times C_2$
нетривиальный $\psi$ определяется своим ядром~--- одной из трёх подгрупп порядка $2$; $\Aut H\cong S_3$ переставляет их
транзитивно, так что группа одна. Если $\ker\psi=\langle z\rangle$ и $s\in H\smallsetminus\langle z\rangle$, то $\langle t,s\rangle\cong S_3$
($sts^{-1}=t^{-1}$), $z$ централен, и $G=\langle t,s\rangle\times\langle z\rangle\cong S_3\times C_2\cong D_6$.
\end{itemize}
Итак, групп порядка $12$ пять: $C_{12}$, $C_2\times C_6$, $A_4$, $C_3\rtimes C_4$, $D_6$ (они попарно неизоморфны: их различают
силовская $2$-подгруппа и числа $n_2,n_3$). Во всех, кроме $A_4$, есть элемент порядка $6$: в абелевых~--- произведение
элементов порядков $2$ и $3$; в $C_3\rtimes C_4$ элемент $s^2$ действует на $t$ тривиально ($\psi(s^2)=\psi(s)^2=1$), и $ts^2$ имеет
порядок $6$; в $C_3\rtimes(C_2\times C_2)$ элемент $z$ коммутирует с $t$, и $tz$ имеет порядок $6$. В $S_4$ элементов порядка $6$ нет
(цикловые типы $1$, $2$, $2+2$, $3$, $4$), поэтому подгруппа порядка $12$ в $S_4$ изоморфна $A_4$. Это согласуется и с (b):
у транзитивной подгруппы $G\subset S_4$ порядка $12$ стабилизатор точки~--- подгруппа порядка $3$, не содержащая неединичных
нормальных подгрупп $G$, т.\,е.\ ненормальная; значит, $n_3\ne1$, $n_3=4$, и $G\cong A_4$ по второму случаю.
\end{remark}

\subsection*{Группа Галуа неприводимой квартики}

\begin{theorem}[4.9d]
Пусть $\operatorname{char}F\ne2$, $f\in F[X]$~--- неприводимый сепарабельный многочлен степени $4$, $g$~--- его кубическая
резольвента, $M$~--- поле разложения $g$, $D=D(f)=D(g)$. Тогда
\begin{center}
\begin{tabular}{@{}lll@{}}
\toprule
$g$ над $F$ & $D$ & $G_f$\\
\midrule
неприводим & не квадрат & $S_4$\\
неприводим & квадрат & $A_4$\\
раскладывается на линейные множители & квадрат & $V$\\
ровно один корень в $F$ & не квадрат & $D_4$, если $f$ неприводим над $M$;\\
 & & $C_4$, если $f$ приводим над $M$\\
\bottomrule
\end{tabular}
\end{center}
В последнем случае $M=F[\sqrt D]$, и $G_f=D_4\iff[E:M]=4$, $G_f=C_4\iff[E:M]=2$.
\end{theorem}

\begin{proof}
По \rf{4.5} $G=G_f$ транзитивна, т.\,е.\ это одна из групп \rf{4.9b}, и по \rf{4.8} $\Gal(M/F)\cong G/(G\cap V)$. Многочлен $g$ сепарабелен
(\rf{4.7j}(a)), а $\Gal(M/F)$ действует на его корнях так же, как $G$ (через $\pi$).

Если $G\in\{S_4,A_4\}$, то $\pi(G)$~--- это $S_3$ или $A_3$, оно транзитивно на $\theta_1,\theta_2,\theta_3$, и $g$ неприводим (\rf{4.5}).
Если $G=V$, то $\pi(G)=1$, т.\,е.\ $M=F$ и $g$ раскладывается на линейные множители. Если $G\in\{D_4,C_4\}$, то $\pi(G)$ имеет
порядок $2$, т.\,е.\ это транспозиция: ровно один корень $g$ неподвижен, т.\,е.\ лежит в $F$ (\rf{3.17}). Эти три случая
исключают друг друга и исчерпывают все возможности (у кубического многочлена не бывает ровно двух корней в $F$), так что
поведение $g$ определяет, в какой из групп $\{S_4,A_4\}$, $\{V\}$, $\{D_4,C_4\}$ лежит $G$.

$S_4$ и $A_4$ различаются условием $G\subset A_4$, т.\,е.\ (\rf{4.2}) тем, является ли $D$ квадратом. Группы $V\subset A_4$ и
$D_4,C_4\not\subset A_4$ дают столбец $D$ в остальных строках.

Пусть $G\in\{D_4,C_4\}$. Тогда $[M:F]=2$ и $\Delta(f)=\Delta(g)\in M$ (\rf{4.9}), $\Delta(f)\notin F$ (так как $G\not\subset A_4$), откуда
$M=F[\Delta(f)]=F[\sqrt D]$. Далее $\Gal(E/M)=G\cap V$, и это $V$ для $D_4$ (транзитивна) и $\langle(1\,2)(3\,4)\rangle$ для $C_4$
(нетранзитивна: орбиты $\{x_1,x_2\}$, $\{x_3,x_4\}$). По \rf{4.5}, применённому над $M$, $f$ неприводим над $M$ в первом случае и
приводим во втором. Наконец, $[E:M]=|G\cap V|$ равно $4$ или $2$.
\end{proof}

Итак, всё определяется одной кубической резольвентой и дискриминантом, кроме выбора между $D_4$ и $C_4$, для которого
нужна дополнительная информация. Её можно получить, не раскладывая $f$ над $M$. Для этого посмотрим на подполя.

\begin{proposition}[4.9e]
Пусть в условиях \rf{4.9d} $g$ имеет ровно один корень в $F$; занумеруем корни $f$ так, что это $\theta_1=x_1x_2+x_3x_4$. Тогда
$G\subset D^{(1)}$. Положим
\[
\delta_1=(x_1x_2-x_3x_4)^2=\theta_1^2-4e,\qquad \delta_2=(x_1+x_2-x_3-x_4)^2=b^2-4(c-\theta_1).
\]
Оба числа лежат в $F$, и хотя бы одно из них не равно $0$. Подгруппа $C_4=\langle(1\,3\,2\,4)\rangle$ оставляет на месте
элементы $\eta_1=\Delta(f)\,(x_1x_2-x_3x_4)$ и $\eta_2=\Delta(f)\,(x_1+x_2-x_3-x_4)$, а элементы $D^{(1)}\smallsetminus C_4$ меняют их знак;
$\eta_j^2=D\delta_j$.
\end{proposition}

\begin{proof}
$G$ оставляет на месте $\theta_1\in F$, а $\theta_i$ различны, поэтому $\pi(G)$ оставляет на месте $1$ и $G\subset D^{(1)}$. Так как
$x_1x_2+x_3x_4=\theta_1$, $x_1x_2\cdot x_3x_4=e$, $(x_1+x_2)+(x_3+x_4)=-b$ и $(x_1+x_2)(x_3+x_4)=c-\theta_1$ (см.\ доказательство
\rf{4.9}), получаем формулы для $\delta_1,\delta_2$. Если $\delta_1=\delta_2=0$, то $x_1x_2=x_3x_4$ и $x_1+x_2=x_3+x_4$, т.\,е.\
$(X-x_1)(X-x_2)=(X-x_3)(X-x_4)$, что невозможно, так как корни различны. Наконец, $(1\,3\,2\,4)$~--- нечётная перестановка
($\Delta\mapsto-\Delta$), переводящая $x_1x_2-x_3x_4$ в $x_3x_4-x_1x_2$ и $x_1+x_2-x_3-x_4$ в противоположное; $(1\,2)(3\,4)$ чётна и
сохраняет обе разности; а $(1\,2)$ нечётна и сохраняет обе разности. Так как $D^{(1)}=C_4\cup(1\,2)C_4$, всё доказано.
\end{proof}

\begin{keepfig}
\begin{center}
\begin{tikzpicture}[every node/.style={font=\small},xscale=1.15]
\node (one) at (0,3.6) {$1$};
\node (a) at (-4,2.2) {$\langle(1\,2)\rangle$};
\node (b) at (-2,2.2) {$\langle(3\,4)\rangle$};
\node (c) at (0,2.2) {$\langle(1\,2)(3\,4)\rangle$};
\node (d) at (2.4,2.2) {$\langle(1\,3)(2\,4)\rangle$};
\node (e) at (4.6,2.2) {$\langle(1\,4)(2\,3)\rangle$};
\node (Vp) at (-3,0.8) {$V'=\langle(1\,2),(3\,4)\rangle$};
\node (C) at (0,0.8) {$C_4=\langle(1\,3\,2\,4)\rangle$};
\node (V) at (4.3,0.8) {$V$};
\node[cold] (D) at (0,-0.6) {$D^{(1)}$};
\foreach \x in {a,b,c,d,e}{\draw (one)--(\x);}
\draw (a)--(Vp); \draw (b)--(Vp); \draw (c)--(Vp); \draw (c)--(C); \draw (c)--(V); \draw (d)--(V); \draw (e)--(V);
\draw (Vp)--(D); \draw (C)--(D); \draw (V)--(D);
\node[cnew,font=\scriptsize,anchor=east] at (Vp.west) {$F[\sqrt{\delta_j}]\leftrightarrow$};
\node[cnew,font=\scriptsize,anchor=west] at (C.east) {$\leftrightarrow F[\sqrt{D\delta_j}]$};
\node[cnew,font=\scriptsize,anchor=west] at (V.east) {$\leftrightarrow F[\sqrt D]=M$};
\end{tikzpicture}
\end{center}
\figcap{Подгруппы группы $D^{(1)}=\langle(1\,3\,2\,4),(1\,2)\rangle\cong D_4$ ($1$ сверху). При $G_f=D^{(1)}$ трём подгруппам порядка $4$ отвечают
три квадратичных подполя поля разложения: $V=G\cap A_4\leftrightarrow F[\sqrt D]$, $V'\leftrightarrow F[x_1x_2-x_3x_4]$ или
$F[x_1+x_2-x_3-x_4]$, $C_4\leftrightarrow F[\eta_j]$; здесь $\delta_j\ne0$ (\rf{4.9e}). Группа $G_f=C_4$~--- это случай, когда $\eta_j\in F$.}
\end{keepfig}

\begin{theorem}[4.9f: критерий Каппе--Уоррена]
Пусть $\operatorname{char}F\ne2$, $f=X^4+bX^3+cX^2+dX+e\in F[X]$ неприводим и сепарабелен, и его кубическая резольвента имеет
ровно один корень $\theta\in F$. Тогда
\[
G_f=C_4\iff\text{оба числа }(\theta^2-4e)\,D\ \text{ и }\ \bigl(b^2-4(c-\theta)\bigr)D\ \text{являются квадратами в }F;
\]
в противном случае $G_f=D_4$.
\end{theorem}

\begin{proof}
В обозначениях \rf{4.9e} ($\theta=\theta_1$) нужно доказать: $G=C_4\iff D\delta_1$ и $D\delta_2$~--- квадраты в $F$. По \rf{4.9d} $G=C_4$ или
$G=D^{(1)}$. Если $G=C_4$, то $\eta_1,\eta_2\in E^G=F$ (\rf{4.9e}), и $D\delta_j=\eta_j^2$~--- квадраты. Если $G=D^{(1)}$, выберем $j$ с $\delta_j\ne0$;
тогда $\eta_j\ne0$, и $(1\,2)\in G$ переводит $\eta_j$ в $-\eta_j\ne\eta_j$, т.\,е.\ $\eta_j\notin F$. Если бы $D\delta_j=\eta_j^2$ равнялось $\rho^2$ с
$\rho\in F$, то $\eta_j=\pm\rho\in F$; значит, $D\delta_j$ не квадрат.
\end{proof}

\begin{remark}[4.9g]
(a) Эквивалентная формулировка: \emph{$G_f=C_4$ тогда и только тогда, когда оба многочлена $X^2-\theta X+e$ и
$X^2+bX+(c-\theta)$ раскладываются на линейные множители над $F[\sqrt D]$.} Их корни~--- $x_1x_2,x_3x_4$ и $x_1+x_2,x_3+x_4$. При
$G=C_4$ они неподвижны относительно $G\cap A_4=\langle(1\,2)(3\,4)\rangle$, т.\,е.\ лежат в $E^{G\cap A_4}=F[\sqrt D]$ (\rf{4.2}). При $G=D_4$ они
не могут все лежать в $F[\sqrt D]=E^V$: иначе они неподвижны относительно $(1\,3)(2\,4)\in V$, и $x_1x_2=x_3x_4$, $x_1+x_2=x_3+x_4$,
что невозможно (\rf{4.9e}).

(b) Другой способ различить $D_4$ и $C_4$~--- найти степень поля разложения: $|G_f|=[E:F]$ равно $8$ или $4$.

(c) Таблица \rf{4.9d} относится только к \emph{неприводимым} $f$. Для приводимого $f$ группа $G_f$ нетранзитивна, и её находят по
множителям (\rf{3.27e}). Например, если $f$~--- произведение двух неприводимых квадратичных многочленов с разными
квадратичными полями, то $G_f=\{1,(1\,2),(3\,4),(1\,2)(3\,4)\}\cong C_2\times C_2$~--- группа, изоморфная $V$, но не равная ей и не
лежащая в $A_4$ (см.\ \rf{4.14}).
\end{remark}

\subsection*{Решение уравнения четвёртой степени}

Теория выше подсказывает, как решать уравнение $f=0$: корень резольвенты $\theta_1$ (находится по формуле Кардано)
выделяет разбиение $\{12|34\}$, т.\,е.\ разложение $f$ на два квадратичных множителя, которые находятся квадратными
корнями. Будем считать $f=X^4+pX^2+qX+r$ приведённым (\rf{4.9a}(b)); его резольвента $g=X^3-pX^2-4rX+(4pr-q^2)$. Положим
\[
z_i=\theta_i-p ;\qquad\text{тогда }z_1=(x_1+x_2)^2=(x_3+x_4)^2,\quad z_2=(x_1+x_3)^2,\quad z_3=(x_1+x_4)^2,
\]
так как $p=\theta_1+\theta_2+\theta_3$ и $\theta_2+\theta_3=(x_1+x_2)(x_3+x_4)=-(x_1+x_2)^2$. Числа $z_i$~--- корни многочлена
\[
h(Z)=g(Z+p)=Z^3+2pZ^2+(p^2-4r)Z-q^2 .
\]

\begin{proposition}[4.9h: метод Феррари]
Пусть $\theta$~--- корень $g$, $y_0=\frac\theta2$ и $m=2y_0-p=\theta-p$. Если $m\ne0$, то
\[
f=\Bigl(X^2-\sqrt m\,X+y_0+\frac q{2\sqrt m}\Bigr)\Bigl(X^2+\sqrt m\,X+y_0-\frac q{2\sqrt m}\Bigr),
\]
и корни $f$~--- корни этих двух квадратных трёхчленов. Если $m=0$ для $\theta=\theta_1$, то $q=0$, и $f=X^4+pX^2+r$ решается как
квадратное уравнение относительно $X^2$.
\end{proposition}

\begin{proof}
Перепишем $f=0$ как $X^4=-pX^2-qX-r$ и прибавим к обеим частям $2yX^2+y^2$:
\[
(X^2+y)^2=(2y-p)X^2-qX+(y^2-r).
\]
Левая часть~--- полный квадрат при любом $y$; правая~--- квадрат линейного многочлена (при $2y\ne p$) тогда и только тогда,
когда её дискриминант равен нулю: $q^2=4(2y-p)(y^2-r)$, т.\,е.\ $8y^3-4py^2-8ry+4pr-q^2=0$. Это в точности $g(2y)=0$. При
$y=y_0$ правая часть равна $m\bigl(X-\frac q{2m}\bigr)^2$, и $(X^2+y_0)^2-m\bigl(X-\frac q{2m}\bigr)^2$~--- произведение указанных
трёхчленов (разность квадратов). Если $m=\theta_1-p=z_1=0$, то $x_1+x_2=0=x_3+x_4$, и $q=-(x_1+x_2)(x_1+x_3)(x_1+x_4)=0$ (см.\
\rf{4.9j} ниже).
\end{proof}

Смысл: при $\theta=\theta_1$ имеем $\sqrt m=\pm(x_1+x_2)$, и трёхчлены~--- это $(X-x_1)(X-x_2)$ и $(X-x_3)(X-x_4)$. Резольвента
появляется потому, что разбиений на пары три, и метод <<не знает>>, какое выбрать; годится любое.

\begin{proposition}[4.9i: метод Декарта; критерий разложимости]
Приведённый многочлен $f=X^4+pX^2+qX+r\in F[X]$ раскладывается в $F[X]$ в произведение двух квадратичных многочленов
тогда и только тогда, когда либо $h(Z)=Z^3+2pZ^2+(p^2-4r)Z-q^2$ имеет корень вида $k^2$ с $k\in F^\times$, либо $q=0$ и
$p^2-4r$~--- квадрат в $F$.
\end{proposition}

\begin{proof}
Пусть $f=(X^2+kX+l)(X^2-kX+m)$ (коэффициенты при $X$ противоположны, так как в $f$ нет $X^3$). Сравнение коэффициентов:
$l+m-k^2=p$, $k(m-l)=q$, $lm=r$. Если $k\ne0$, то $l+m=p+k^2$, $m-l=\frac qk$, и $4lm=(p+k^2)^2-\frac{q^2}{k^2}=4r$; умножив на $k^2$,
получаем $k^6+2pk^4+(p^2-4r)k^2-q^2=0$, т.\,е.\ $h(k^2)=0$. Если $k=0$, то $q=0$, а $l,m$~--- корни $T^2-pT+r$, так что
$p^2-4r=(l-m)^2$. Обратно, если $h(k^2)=0$, $k\in F^\times$, то $l=\frac12\bigl(p+k^2-\frac qk\bigr)$, $m=\frac12\bigl(p+k^2+\frac qk\bigr)$ дают
разложение (равенство $lm=r$ равносильно $h(k^2)=0$); если $q=0$ и $p^2-4r=s^2$, то $f=\bigl(X^2+\frac{p+s}2\bigr)\bigl(X^2+\frac{p-s}2\bigr)$.
\end{proof}

Вместе с поиском корней это даёт быстрый способ доказывать неприводимость квартик над $\Q$ (примеры \rf{4.12}, \rf{4.13a}).

\begin{theorem}[4.9j: формула Эйлера]
Пусть $z_1,z_2,z_3$~--- корни $h$, и квадратные корни $\sqrt{z_i}$ выбраны так, что $\sqrt{z_1}\sqrt{z_2}\sqrt{z_3}=-q$. Тогда корни $f$~--- это
\[
\tfrac12\bigl(\varepsilon_1\sqrt{z_1}+\varepsilon_2\sqrt{z_2}+\varepsilon_3\sqrt{z_3}\bigr),\qquad \varepsilon_i=\pm1,\quad \varepsilon_1\varepsilon_2\varepsilon_3=1 .
\]
\end{theorem}

\begin{proof}
Возьмём $s_1=x_1+x_2$, $s_2=x_1+x_3$, $s_3=x_1+x_4$; тогда $s_i^2=z_i$ и
\[
s_1s_2s_3=x_1^3+x_1^2(x_2+x_3+x_4)+x_1(x_2x_3+x_2x_4+x_3x_4)+x_2x_3x_4=x_1\cdot(x_1^2+p)+(-q-x_1(x_1^2+p))=-q,
\]
так как $x_2+x_3+x_4=-x_1$, $x_2x_3+x_2x_4+x_3x_4=p+x_1^2$ и $x_2x_3x_4=-q-x_1(x_2x_3+x_2x_4+x_3x_4)$. Далее
$s_1+s_2+s_3=2x_1+(x_1+x_2+x_3+x_4)=2x_1$, $s_1-s_2-s_3=2x_2$, $-s_1+s_2-s_3=2x_3$, $-s_1-s_2+s_3=2x_4$. Любой другой выбор
квадратных корней с произведением $-q$ отличается от $(s_1,s_2,s_3)$ знаками с произведением $1$ (если $q\ne0$; при $q=0$ один
из $s_i$ равен нулю, и знак при нём роли не играет) и лишь переставляет четыре числа.
\end{proof}

\begin{remark}[4.9k: башня решения и цепочка подгрупп]
Числа $2s_1=x_1+x_2-x_3-x_4$, $2s_2$, $2s_3$~--- аналоги резольвент Лагранжа (\rf{4.7d}) для группы $V\cong C_2\times C_2$: каждый элемент
$V$ меняет знак у двух из них или ни у одного, так что их квадраты $4z_i$ неподвижны относительно $V$ и лежат в $M=E^{G\cap V}$.
Решение следует цепочке $S_4\supset A_4\supset V\supset\langle(1\,2)(3\,4)\rangle\supset1$ с факторами $C_2,C_3,C_2,C_2$:
\begin{center}
\begin{tabular}{@{}lll@{}}
\toprule
шаг & что присоединяется & поле\\
\midrule
$S_4\supset A_4$ & $\sqrt D$ & $F[\sqrt D]$\\
$A_4\supset V$ & корни резольвенты (кубический корень по Кардано) & $M=F[\theta_1,\theta_2,\theta_3]$\\
$V\supset C_2\supset1$ & $\sqrt{z_1}$, затем $\sqrt{z_2}$ & $E$\\
\bottomrule
\end{tabular}
\end{center}
(над полем, содержащим $\omega$). Разрешимость $S_4$ и делает возможной формулу в радикалах (\rf{3.28}).
\end{remark}

\begin{aside}[4.9l: пучок коник]
Положим $s=x^2$. Корни $f$~--- абсциссы четырёх точек $P_i=(x_i,x_i^2)$ пересечения параболы $s-x^2=0$ с коникой
$s^2+ps+qx+r=0$. Через эти точки проходят все коники пучка $C_\lambda\colon s^2+ps+qx+r+\lambda(x^2-s)=0$. Коника $C_\lambda$ вырождается в
пару прямых, когда
\[
\det\begin{pmatrix}\lambda&0&q/2\\0&1&(p-\lambda)/2\\q/2&(p-\lambda)/2&r\end{pmatrix}=-\tfrac14\bigl(\lambda^3-2p\lambda^2+(p^2-4r)\lambda+q^2\bigr)=\tfrac14h(-\lambda)=0,
\]
т.\,е.\ при $\lambda=-z_i$. Три вырожденные коники~--- пары прямых $P_1P_2\cup P_3P_4$, $P_1P_3\cup P_2P_4$, $P_1P_4\cup P_2P_3$, по одной на
разбиение на пары. Метод Феррари~--- это нахождение вырожденной коники (кубическое уравнение), её разложение на две прямые
(квадратный корень) и пересечение каждой прямой с параболой (квадратные уравнения).
\end{aside}

\begin{example}[4.9m]
$f=X^4-2X^2+8X-3$: $p=-2$, $q=8$, $r=-3$. Уравнение $g(2y)=0$ из \rf{4.9h}: $8y^3+8y^2+24y-40=8(y-1)(y^2+2y+5)=0$, корень $y_0=1$, $m=2y_0-p=4$,
$\sqrt m=2$, и
\[
(X^2+1)^2=4X^2-8X+4=(2X-2)^2\ \Longrightarrow\ f=(X^2-2X+3)(X^2+2X-1),
\]
корни $1\pm i\sqrt2$, $-1\pm\sqrt2$.
\end{example}

\begin{proposition}[4.9n: вещественные корни]
Пусть $p,q,r\in\R$ и $D=D(X^4+pX^2+qX+r)\ne0$. Тогда
\begin{itemize}
\item $D<0$ $\iff$ у $f$ два вещественных и два невещественных корня;
\item $D>0$, $p<0$ и $p^2-4r>0$ $\iff$ все четыре корня вещественны;
\item $D>0$ и ($p\ge0$ или $p^2-4r\le0$) $\iff$ вещественных корней нет.
\end{itemize}
\end{proposition}

\begin{proof}
Первое утверждение и то, что при $D>0$ число невещественных корней равно $0$ или $4$,~--- это \rf{4.3}. Пусть $D>0$. Тогда $D(h)=D(g)=D>0$,
и $z_1,z_2,z_3$~--- различные вещественные числа. Если все $x_i$ вещественны, то $z_i=(x_1+x_{i+1})^2\ge0$, причём нулю равно не более
одного $z_i$; значит, $2p=-\sum z_i<0$ и $p^2-4r=\sum_{i<j}z_iz_j>0$. Обратно, пусть $p<0$, $p^2-4r>0$, но корни~--- $w,\bar w,v,\bar v$ с
$w,v\notin\R$. Для разбиения $\{w,v|\bar w,\bar v\}$ имеем $w+v=-(\bar w+\bar v)$, т.\,е.\ $w+v$ чисто мнимо и соответствующее $z=(w+v)^2\le0$. Но
при $z<0$ все слагаемые $h(z)=z^3+2pz^2+(p^2-4r)z-q^2$ отрицательны или равны нулю, а первое отрицательно; поэтому $z=0$, $v=-w$, и
$f=(X^2-w^2)(X^2-\bar w^2)$, $p=-2\operatorname{Re}(w^2)$, $r=|w|^4$, $p^2-4r=-4(\operatorname{Im}w^2)^2\le0$~--- противоречие.
\end{proof}

\subsection*{Примеры}

\begin{example}[4.10]
$f=X^4-4X+2\in\Q[X]$ неприводим по критерию Эйзенштейна (\rf{1.16}, $p=2$). По \rf{4.9} его резольвента
$g=X^3-8X-16$; она неприводима, так как не имеет корней даже по модулю $5$ (её значения в $0,1,2,3,4$ по модулю $5$ равны
$4,2,1,2,1$). Дискриминант $D(g)=-4(-8)^3-27\cdot16^2=-4864=-2^8\cdot19$ не квадрат. По \rf{4.9d} $G_f=S_4$ (ср.\ \rf{3.27}, где отсюда
выведено, что корни $f$ не построимы).
\end{example}

\begin{example}[4.11]
$f=X^4+4X^2+2\in\Q[X]$ неприводим по критерию Эйзенштейна. Его резольвента $X^3-4X^2-8X+32=(X-4)(X^2-8)$ имеет ровно
один рациональный корень $\theta=4$, и $M=\Q[\sqrt8]=\Q[\sqrt2]$; $D=16\cdot2\cdot(16-8)^2=2^{11}$ (\rf{4.2f}(a)). Над $M$
\[
f=(X^2+2-\sqrt2)(X^2+2+\sqrt2)
\]
(уже как многочлен от $X^2$: $X^2=-2\pm\sqrt2$), так что $G_f=C_4$. Проверка по \rf{4.9f}: $(\theta^2-4e)D=8\cdot2^{11}=2^{14}=(2^7)^2$ и
$b^2-4(c-\theta)=0$. Явно: корни $\pm\alpha$, $\pm\alpha'$, где $\alpha=i\sqrt{2+\sqrt2}$, $\alpha'=i\sqrt{2-\sqrt2}$; так как $\sqrt2=-\alpha^2-2\in\Q[\alpha]$ и
$\alpha\alpha'=-\sqrt2$, все корни лежат в $\Q[\alpha]$, т.\,е.\ $E=\Q[\alpha]$ степени $4$. Автоморфизм $\sigma\colon\alpha\mapsto\alpha'$ переводит $\sqrt2$ в
$-\sigma(\alpha)^2-2=-\sqrt2$, а значит, $\sigma^2(\alpha)=\sigma\bigl(-\frac{\sqrt2}\alpha\bigr)=\frac{\sqrt2}{\alpha'}=-\alpha$: $\sigma$ имеет порядок $4$.
\end{example}

\begin{example}[4.11a]
Аналогично $X^4-4X^2+2$ (минимальный многочлен числа $2\cos\frac\pi8=\sqrt{2+\sqrt2}$, корни $\pm\sqrt{2\pm\sqrt2}$) имеет резольвенту
$(X+4)(X^2-8)$, $D=2^{11}$, $\theta=-4$, $(\theta^2-4e)D=2^{14}$, и $G=C_4$. Поле $\Q[2\cos\frac\pi8]$ вещественно, и корни получаются
<<вложенными>> квадратными корнями $\Q\subset\Q[\sqrt2]\subset\Q[\sqrt{2+\sqrt2}]$: сопряжённый корень $\sqrt{2-\sqrt2}=\sqrt2/\sqrt{2+\sqrt2}$ нового
радикала не требует.
\end{example}

\begin{example}[4.12]
$f=X^4-10X^2+4\in\Q[X]$. \emph{Неприводимость.} Рациональные корни были бы целыми делителями $4$, но
$f(\pm1)=-5$, $f(\pm2)=-20$, $f(\pm4)=100$. По \rf{4.9i} ($p=-10$, $q=0$, $r=4$) остаётся проверить, что $p^2-4r=84$ не квадрат и что
$h(Z)=Z^3-20Z^2+84Z=Z(Z-6)(Z-14)$ не имеет ненулевых корней-квадратов: $6$ и $14$~--- не квадраты в $\Q$. Значит, $f$ неприводим.

\emph{Группа.} Резольвента $X^3+10X^2-16X-160=(X+10)(X+4)(X-4)$ раскладывается над $\Q$, и по \rf{4.9d} $G_f=V$. Явно: $X^2=5\pm\sqrt{21}$,
корни~--- $x_{1,2}=\pm\sqrt{5+\sqrt{21}}$, $x_{3,4}=\pm\sqrt{5-\sqrt{21}}$, и $\theta_1=-(5+\sqrt{21})-(5-\sqrt{21})=-10$,
$\theta_{2,3}=\pm2\sqrt{(5+\sqrt{21})(5-\sqrt{21})}=\pm4$. Так как $\sqrt{5\pm\sqrt{21}}=\frac{\sqrt{14}\pm\sqrt6}2$, поле разложения~--- $\Q[\sqrt6,\sqrt{14}]$,
биквадратичное поле; $D=16\cdot4\cdot84^2=672^2$~--- квадрат, как и должно быть.
\end{example}

\begin{example}[4.13]
$f=X^4-2\in\Q[X]$ неприводим по критерию Эйзенштейна; резольвента $g=X^3+8X=X(X^2+8)$ имеет ровно один рациональный
корень $\theta=0$, и $M=\Q[\sqrt{-8}]=\Q[\sqrt{-2}]$, $D=-2^{11}$. По \rf{4.9d} $G_f=D_4$ или $C_4$.

\emph{$f$ неприводим над $M$.} Корни $f$~--- $\pm\alpha,\pm i\alpha$, $\alpha=\sqrt[4]2$. Для любого корня $\beta$ поле $M[\beta]$ содержит $\beta^2=\pm\sqrt2$,
затем $\frac{\sqrt{-2}}{\sqrt2}=i$ и $\alpha\in\{\pm\beta,\pm i\beta\}$, т.\,е.\ $M[\beta]=\Q[\alpha,i]$. Степень этого поля равна $8$: $[\Q[\alpha]:\Q]=4$ и
$i\notin\Q[\alpha]\subset\R$. Поэтому $[M[\beta]:M]=4$, и минимальный многочлен любого корня над $M$ имеет степень $4$, т.\,е.\ равен $f$.
Значит, $G_f=D_4$. (Короче: $[E:\Q]=[\Q[\alpha,i]:\Q]=8$, а транзитивная подгруппа $S_4$ порядка $8$~--- это $D_4$.)

\emph{Проверка по \rf{4.9f}:} $(\theta^2-4e)D=8\cdot(-2^{11})=-2^{14}$ не квадрат. Три квадратичных подполя (рис.\ к \rf{4.9e}):
$\Q[\sqrt D]=\Q[\sqrt{-2}]$, $\Q[\sqrt{\delta_1}]=\Q[\sqrt8]=\Q[\sqrt2]$, $\Q[\sqrt{D\delta_1}]=\Q[\sqrt{-2^{14}}]=\Q[i]$.
\end{example}

\begin{example}[4.13a]
$f=X^4+8X+12\in\Q[X]$. Так как $f=(X^2-2)^2+4(X+1)^2+4>0$ на $\R$, у $f$ нет вещественных (тем более рациональных) корней.
Здесь $p=0$, и $h=g=X^3-48X-64$; у $g$ нет рациональных корней (среди делителей $64$ корней нет; короче: по модулю $5$ имеем
$g\equiv X^3+2X+1$, и его значения в $0,1,2,3,4$ равны $1,4,3,4,3$). По \rf{4.9i} $f$ не
раскладывается на квадратичные множители, т.\,е.\ неприводим. Резольвента неприводима, $D(g)=-4(-48)^3-27\cdot64^2=331776=576^2$~---
квадрат. По \rf{4.9d} $G_f=A_4$.
\end{example}

\begin{example}[4.13b]
Круговой многочлен $\Phi_5=X^4+X^3+X^2+X+1$: резольвента $X^3-X^2-3X+2=(X-2)(X^2+X-1)$, $\theta=2$, $D=125=5^3$. По \rf{4.9f}:
$\theta^2-4e=0$ и $\bigl(b^2-4(c-\theta)\bigr)D=5\cdot125=25^2$; значит, $G=C_4$~--- в согласии с $\Gal(\Q[\zeta_5]/\Q)\cong(\Z/5\Z)^\times$. Над
$\Q[\sqrt D]=\Q[\sqrt5]$ имеем $\Phi_5=\bigl(X^2+\frac{1-\sqrt5}2X+1\bigr)\bigl(X^2+\frac{1+\sqrt5}2X+1\bigr)$. Здесь $b\ne0$, и полезен именно второй
многочлен из \rf{4.9g}(a).
\end{example}

\begin{remark}[4.13c: сводка]
Для $X^4-X-1$ (резольвента $X^3+4X-1$ неприводима, $D=-283$) группа равна $S_4$ (\rf{3.27}). Как группа определяет способ
решения (\rf{4.9k}):
\begin{center}
\begin{tabular}{@{}llp{7.2cm}@{}}
\toprule
$G$ & пример & что нужно для корней\\
\midrule
$V$ & $X^4-10X^2+4$ (\rf{4.12}) & два независимых квадратных корня\\
$C_4$ & $X^4+4X^2+2$, $\Phi_5$ & вложенные квадратные корни\\
$D_4$ & $X^4-2$ (\rf{4.13}) & вложенные и независимые квадратные корни\\
$A_4$ & $X^4+8X+12$ & циклическая кубика ($D(g)$~--- квадрат) и квадратные корни\\
$S_4$ & $X^4-4X+2$, $X^4-X-1$ & $\sqrt D$, кубическая резольвента, квадратные корни\\
\bottomrule
\end{tabular}
\end{center}
Корни построимы циркулем и линейкой ровно в первых трёх случаях (\rf{3.25b}: $|G|$~--- степень двойки).
\end{remark}

\subsection*{Задача о скрещённых лестницах}

\begin{example}[4.14]
Пусть $c,d\in\Q$, $c>0$, $d>0$, и
\[
f(X)=X^4-2cX^3-dX^2+2cdX-dc^2 .
\]
Положим $r=d/c^2$, и пусть $w$~--- единственное положительное вещественное число с $r=\frac{w^3}{w^2+4}$. Пусть $m$~--- число
рациональных корней $f$ (если $c,d\in\Z$, то $f\in\Z[X]$ унитарен, и рациональные корни~--- целые). Тогда
\begin{itemize}
\item если $m=0$ и $w\notin\Q$, то $G_f=S_4$;
\item если $m=1$, то $w\notin\Q$ и $G_f=S_3$;
\item если $w\in\Q$ и $w^2+4$ не квадрат в $\Q$, то $m=0$ и $G_f=D_4$;
\item если $w\in\Q$ и $w^2+4$~--- квадрат в $\Q$, то $m=0$ и $G_f\cong C_2\times C_2$ (нетранзитивная группа
$\{1,(1\,2),(3\,4),(1\,2)(3\,4)\}$).
\end{itemize}
Других случаев нет: $m\ge2$ не бывает.
\end{example}

\begin{remark}[4.14a: откуда этот многочлен]
В переулке ширины $\ell$ стоят две лестницы длин $a>b$; нижний конец каждой упирается в основание одной стены, верхний~---
в противоположную стену на высоте $h_1$ (лестница $a$) и $h_2$ (лестница $b$). Лестницы пересекаются на высоте $c$. Известны
$a,b,c$; найти $h_1$ (а с ним $\ell=\sqrt{a^2-h_1^2}$).

Точка пересечения на расстоянии $u$ от левой стены лежит на обеих лестницах: $c=\frac{h_1u}\ell=\frac{h_2(\ell-u)}\ell$. Отсюда
$\frac c{h_1}+\frac c{h_2}=\frac u\ell+\frac{\ell-u}\ell=1$, т.\,е.\ $\frac1c=\frac1{h_1}+\frac1{h_2}$. Кроме того, $h_1^2-h_2^2=(a^2-\ell^2)-(b^2-\ell^2)=a^2-b^2=:d$.
Подставив $h_2=\frac{ch_1}{h_1-c}$ в $h_1^2-h_2^2=d$ и умножив на $(h_1-c)^2$, получаем $h_1^2(h_1-c)^2-c^2h_1^2=d(h_1-c)^2$, т.\,е.\ ровно
$f(h_1)=0$.

Функция $x\mapsto x^2-\bigl(\frac{cx}{x-c}\bigr)^2$ возрастает на $(c,\infty)$ от $-\infty$ до $+\infty$ (оба слагаемых возрастают) и равна $0$ при $x=2c$;
поэтому $f$ имеет ровно один корень в $(c,\infty)$, и он больше $2c$~--- это $h_1$. На $(0,c)$ эта функция отрицательна, так что второй
вещественный корень $f$ (см.\ ниже) отрицателен. Обратно, при любых $c,d>0$ корень $h_1>2c$ даёт $h_2=\frac{ch_1}{h_1-c}\in(c,h_1)$, и
при любом $a>h_1$ получается настоящая конфигурация лестниц с $\ell=\sqrt{a^2-h_1^2}$, $b=\sqrt{\ell^2+h_2^2}$.
\end{remark}

\begin{keepfig}
\begin{center}
\begin{tikzpicture}[scale=0.85,every node/.style={font=\small}]
\fill[gray!15] (-0.25,0) rectangle (0,5.6); \fill[gray!15] (2.8,0) rectangle (3.05,5.6);
\draw[thick] (0,0)--(0,5.6); \draw[thick] (2.8,0)--(2.8,5.6); \draw[thick] (-0.6,0)--(3.4,0);
\draw[cold,very thick] (0,0)--node[pos=0.78,left]{$a$} (2.8,5.25);
\draw[cold,very thick] (2.8,0)--node[pos=0.3,above]{$b$} (0,2.1);
\fill[cnew] (0.8,1.5) circle (1.6pt);
\draw[cnew,dashed] (0.8,1.5)--(0.8,0) node[midway,right]{$c$};
\draw[<->] (3.3,0)--node[right]{$h_1$} (3.3,5.25); \draw[dotted] (2.8,5.25)--(3.3,5.25);
\draw[<->] (-0.5,0)--node[left]{$h_2$} (-0.5,2.1); \draw[dotted] (0,2.1)--(-0.5,2.1);
\draw[<->] (0,-0.35)--node[below]{$\ell$} (2.8,-0.35);
\draw[<->] (0,-0.9)--node[below]{$u$} (0.8,-0.9);
\end{tikzpicture}
\end{center}
\figcap{Скрещённые лестницы (в масштабе: $a=119$, $b=70$, $c=30$, $\ell=56$, $h_1=105$, $h_2=42$). Данные~--- длины $a,b$ и высота
пересечения $c$; ищется $h_1$. Из подобия $\frac1c=\frac1{h_1}+\frac1{h_2}$, а $h_1^2-h_2^2=a^2-b^2$.}
\end{keepfig}

\begin{proof}[Доказательство 4.14]
\emph{Шаг 1: нормировка.} Подстановка $X=cY$ даёт $f(cY)=c^4g(Y)$, где
\[
g(Y)=Y^4-2Y^3-rY^2+2rY-r,\qquad r=\frac d{c^2}>0 .
\]
Корни $f$ и $g$ отличаются множителем $c\in\Q^\times$, поэтому у них одни и те же поле разложения, группа Галуа и число
рациональных корней; $D(f)=c^{12}D(g)$.

\emph{Шаг 2: резольвента и дискриминант.} По \rf{4.9} (коэффициенты $-2,\,-r,\,2r,\,-r$ многочлена $g$ стоят на местах $b,c,d,e$) резольвента многочлена $g$ равна
\[
R(Z)=Z^3+rZ^2+4r .
\]
Подстановка $Z=-w$ даёт $-w^3+r(w^2+4)$, т.\,е.\ \emph{$-w$~--- корень $R$}; это и есть смысл числа $w$ (функция $w\mapsto\frac{w^3}{w^2+4}$
строго возрастает на $(0,\infty)$ от $0$ до $\infty$, так что $w$ определено однозначно). Далее $D(g)=D(R)=-16r^2(r^2+27)<0$. Значит
(\rf{4.3}), у $g$ ровно два вещественных корня, у $R$~--- один вещественный корень $-w$ и два невещественных. В частности,
\emph{$R$ имеет рациональный корень тогда и только тогда, когда $w\in\Q$, и тогда этот корень единственный}. Кроме того, $g$
сепарабелен, $D(g)$ не квадрат, т.\,е.\ $G\not\subset A_4$, и $m\le2$. Подставив $r=\frac{w^3}{w^2+4}$, получаем (прямое вычисление)
\[
D(g)=-\frac{16\,w^6(w^2+3)(w^2+12)^2}{(w^2+4)^4},\qquad w^2+4r=\frac{w^2(w+2)^2}{w^2+4} .
\]

\emph{Шаг 3: рациональные квадратичные множители.} Пусть $g=(Y^2-sY+P)(Y^2-s'Y+P')$ над $\Q$, т.\,е.\ $\{y_1,y_2\}$, $\{y_3,y_4\}$~--- корни
множителей. Тогда $\theta_1=P+P'\in\Q$~--- рациональный корень $R$, т.\,е.\ $w\in\Q$ и $\theta_1=-w$; при этом $P,P'$~--- корни
$T^2+wT-r$ (так как $PP'=-r$), и дискриминант $w^2+4r=\frac{w^2(w+2)^2}{w^2+4}$ этого трёхчлена должен быть квадратом, т.\,е.\ $w^2+4$~---
квадрат в $\Q$. Обратно, если $w\in\Q$ и $w^2+4=v^2$, $v\in\Q_{>0}$, положим $\lambda=\frac{w+v}2$; тогда $\lambda\in\Q$, $\lambda>1$, $w=\lambda-\lambda^{-1}$,
и прямое вычисление даёт
\[
\lambda(\lambda^2+1)^2\,g=q_1q_2,\qquad
\begin{aligned}
q_1&=(\lambda^2+1)Y^2-2(\lambda+1)Y-(\lambda+1)^2(\lambda-1),\\
q_2&=\lambda(\lambda^2+1)Y^2-2\lambda^2(\lambda-1)Y+(\lambda-1)^2(\lambda+1),
\end{aligned}
\]
причём $D(q_1)=4\lambda(\lambda+1)^2(\lambda^2-\lambda+1)$ и $D(q_2)=-4\lambda(\lambda-1)^2(\lambda^2+\lambda+1)<0$.

\emph{Шаг 4: $m\le1$.} Если $m=2$, то $(Y-k_1)(Y-k_2)$~--- рациональный множитель $g$, и по шагу 3 $w\in\Q$, $w^2+4$~--- квадрат, $g$
пропорционален $q_1q_2$. У $q_2$ вещественных корней нет, поэтому $k_1,k_2$~--- корни $q_1$, и $D(q_1)$~--- квадрат в $\Q$, т.\,е.\
$\lambda(\lambda^2-\lambda+1)$~--- квадрат в $\Q$ (множитель $4(\lambda+1)^2$~--- ненулевой квадрат). По \rf{4.14b} при рациональном $\lambda>0$
это возможно только для $\lambda=1$~--- противоречие с $\lambda>1$.

\emph{Шаг 5: разбор случаев.}

\emph{$m=1$.} Тогда $g=(Y-k)q$, где кубический многочлен $q$ неприводим (иначе $m\ge2$). Группа $G=G_q$ действует на корнях
$q$ как подгруппа $S_3$, и $D(g)=D(q)\,q(k)^2$; так как $D(g)$ не квадрат, $D(q)$ тоже не квадрат, и $G=S_3$ (\rf{4.7}). Если бы $w\in\Q$,
то корень $\theta=-w$ резольвенты, скажем $\theta=y_1y_4+y_2y_3$ при $y_4=k$, был бы неподвижен относительно $G$, т.\,е.\ $G$ сохраняла
бы разбиение $\{14|23\}$; но $S_3$ на корнях $y_1,y_2,y_3$ переводит его в $\{24|13\}$ (транспозицией $(1\,2)$). Значит, $w\notin\Q$.

\emph{$m=0$, $w\notin\Q$.} Линейных множителей нет, а рациональный квадратичный множитель дал бы $w\in\Q$ (шаг 3). Значит, $g$
неприводим. Резольвента $R$ неприводима (кубический многочлен без рациональных корней), $D$ не квадрат, и по \rf{4.9d} $G=S_4$.

\emph{$w\in\Q$, $w^2+4$ не квадрат.} Случай $m=1$ невозможен (он даёт $w\notin\Q$), $m\ge2$~--- тоже (шаг 4), а рационального
квадратичного множителя нет по шагу 3. Значит, $g$ неприводим и $m=0$. Резольвента имеет ровно один рациональный корень $\theta=-w$, и $G\in\{D_4,C_4\}$. По \rf{4.9f}
(здесь $e=-r$):
\begin{multline*}
(\theta^2-4e)\,D(g)=(w^2+4r)\,D(g)=-\frac{16\,w^8(w+2)^2(w^2+12)^2(w^2+3)}{(w^2+4)^5}\\
=-(w^2+3)(w^2+4)\cdot\Bigl(\frac{4w^4(w+2)(w^2+12)}{(w^2+4)^3}\Bigr)^2 ;
\end{multline*}
это отрицательное число, и потому не квадрат. Значит, $G=D_4$.

\emph{$w\in\Q$, $w^2+4$~--- квадрат.} По шагам 3 и 4 $g$ пропорционален $q_1q_2$, оба множителя неприводимы ($m=0$), их поля
разложения $\Q[\sqrt{\lambda(\lambda^2-\lambda+1)}]\subset\R$ и $\Q[\sqrt{-\lambda(\lambda^2+\lambda+1)}]\not\subset\R$ различны. Поле разложения $g$~---
их композит степени $4$, и $G=\{1,(1\,2),(3\,4),(1\,2)(3\,4)\}$, где $\{y_1,y_2\}$ и $\{y_3,y_4\}$~--- корни $q_1$ и $q_2$ (\rf{3.27e}).

Так как $m\in\{0,1\}$ и при $m=1$ обязательно $w\notin\Q$, перечисленные случаи исчерпывают все возможности.
\end{proof}

\begin{aside}[4.14b: арифметика шага 4]
Шаг 4~--- единственное место, где нужна не теория Галуа, а арифметика. Там использовано следующее утверждение:
\emph{если $\lambda\in\Q$, $\lambda>0$ и $\lambda(\lambda^2-\lambda+1)$~--- квадрат в $\Q$, то $\lambda=1$.}

Сведём его к классическому диофантову уравнению. Пусть $\lambda=\frac st$, $s,t\in\Z_{>0}$, $\gcd(s,t)=1$. Тогда
$\lambda(\lambda^2-\lambda+1)=\frac{st(s^2-st+t^2)}{t^4}$~--- квадрат в $\Q$ тогда и только тогда, когда натуральное число $st(s^2-st+t^2)$~---
квадрат рационального, а значит (\rf{1.11}), и целого числа. Три положительных сомножителя попарно взаимно просты:
$\gcd(s,s^2-st+t^2)=\gcd(s,t^2)=1$, $\gcd(t,s^2-st+t^2)=\gcd(t,s^2)=1$. По однозначности разложения на простые множители каждый
из них~--- квадрат: $s=x^2$, $t=y^2$, $s^2-st+t^2=z^2$ с натуральными $x,y,z$, т.\,е.
\[
x^4-x^2y^2+y^4=z^2,\qquad \gcd(x,y)=1 .
\]
Дальше мы используем \emph{без доказательства} классическую теорему теории чисел: \emph{в натуральных числах с $\gcd(x,y)=1$
это уравнение имеет единственное решение $x=y=z=1$}. Она доказывается бесконечным спуском (как теорема Ферма о
неразрешимости $x^4+y^4=z^2$); доказательство элементарно, но длинно и к теории Галуа отношения не имеет. Из неё $s=t=1$,
т.\,е.\ $\lambda=1$.

На языке эллиптических кривых: решение $\lambda>0$ с $\lambda(\lambda^2-\lambda+1)=\mu^2$~--- это рациональная точка $(\lambda,\mu)$ кривой
$E\colon y^2=x^3-x^2+x$ (у всех точек $E$ абсцисса $x\ge0$, так как $x^2-x+1>0$). Сформулированная теорема равносильна
равенству $E(\Q)=\{\infty,(1,1),(0,0),(1,-1)\}$: из теоремы оно следует по доказанному выше, а обратно, решению $(x,y,z)$
отвечает точка $\bigl(\frac{x^2}{y^2},\frac{xz}{y^3}\bigr)\in E(\Q)$, и равенство даёт $x^2=y^2$, т.\,е.\ $x=y=1$. Эта группа
циклическая порядка $4$ с образующей $(1,1)$ (касательная $y=x$ в точке $(1,1)$ пересекает $E$ ещё в $(0,0)$, так что
$2\cdot(1,1)=(0,0)$); то, что других рациональных точек нет (ранг $E(\Q)$ равен $0$), доказывается также стандартным $2$-спуском
через изогению $E\to E'\colon y^2=x^3+2x^2-3x$. Таким образом, вопрос <<бывает ли у $f$ два рациональных корня>> сводится к
диофантову уравнению~--- типичный случай, когда вычисление группы Галуа упирается в арифметику.
\end{aside}

\begin{remark}[4.14c: построимость и примеры]
Высота $h_1$~--- корень $f$ в $(2c,\infty)$. По \rf{3.25b} она построима циркулем и линейкой тогда и только тогда, когда поле разложения
её минимального многочлена имеет степень $2^k$:
\begin{itemize}
\item $w\in\Q$ (группы $D_4$ и $C_2\times C_2$): $h_1$ построима (группы порядков $8$ и $4$; во втором случае $h_1$~--- корень квадратного
трёхчлена $q_1$ над $\Q$, растянутого в $c$ раз);
\item $m=0$, $w\notin\Q$ (группа $S_4$): $h_1$ не построима ($|S_4|=24$);
\item $m=1$: $h_1$ построима тогда и только тогда, когда $h_1$~--- это и есть рациональный корень; иначе $h_1$~--- корень неприводимого
кубического многочлена и не построима (\rf{1.38}).
\end{itemize}
Итак, задача о лестницах решается циркулем и линейкой $\iff$ $w\in\Q$ или $h_1\in\Q$. Примеры:
\begin{itemize}
\item $a=119$, $b=70$, $c=30$: $d=9261$, $f=(X-105)(X^3+45X^2-4536X+79380)$, $m=1$, $G=S_3$, и $h_1=105$ рационален ($\ell=56$, $h_2=42$;
рисунок выше);
\item $c=2$, $d=3$: $f=(X+2)(X^3-6X^2+9X-6)$, $m=1$, $G=S_3$, но рациональный корень отрицателен, а $h_1\approx4{,}196$~--- корень
неприводимого кубического множителя: не построима;
\item $a=30$, $b=20$, $c=8$: $d=500$, $w\approx8{,}27\notin\Q$, $m=0$, $G=S_4$;
\item $c=5$, $d=5$: $r=\frac15$, $w=1$, $w^2+4=5$, $G=D_4$; $c=d=1$: $w=2$, $w^2+4=8$, $G=D_4$;
\item $c=10$, $d=54$: $r=\frac{27}{50}$, $w=\frac32$, $w^2+4=\frac{25}4$, $\lambda=2$, и $f=(X^2-12X-180)(X^2-8X+30)$, $G\cong C_2\times C_2$, $h_1=6+6\sqrt6$.
\end{itemize}
\end{remark}

\begin{remark}[4.14d: чем полезен этот пример]
В одном семействе встречаются все приёмы раздела: кубическая резольвента и её рациональный корень, дискриминант (здесь
всегда отрицательный), различение $D_4$ и $C_4$ (\rf{4.9f}), а также приводимые случаи с нетранзитивными группами $S_3$ и
$C_2\times C_2$. Хорошо видно, что дают корни резольвенты: рациональный корень $-w$~--- это разложение $g$ на два квадратичных
множителя над $\Q[\sqrt{w^2+4}]$ (шаг 3), а над $\Q$~--- лишь когда $w^2+4$ квадрат. Пример предостерегает: группа, изоморфная $V$,
может возникнуть и при дискриминанте-неквадрате~--- но тогда она не равна $V$ и нетранзитивна.
\end{remark}
